\title{3.5-meter Segmented-Mirror Robotic Space Telescope}
\subtitle{Mission White Paper: I. Overall Architecture and Scientific Mission}
\newcommand{\wpauthors}{Yong-Woo Kang, Sang Hyun Lee, Juhan Kim, Jeong-Yeol Han, Sungwook E. Hong, Bongkon Moon, Donguk Song, Juhyung Kang, Myeong-Gu Park, Sang Chul Kim, Chung-Uk Lee, Sangmo Tony Sohn, Arman Shafieloo, David Parkinson, Hong Soo Park, Dohyeong Kim, Chan Park, Jungjoo Sohn, Young-Beom Jeon, Jong-Hak Woo, Hyung Mok Lee, Hong Bae Ann, Myungkook James Jee, Mansoo Choi, Changbom Park
}
\date{2026}
%% --- abstract + keywords (use after \maketitle) ---

\newcommand{\wpabstract}{%
We present the preliminary science concept and mission architecture of a 3.5-meter segmented-mirror robotic space telescope currently under study. The telescope is conceived as a versatile space platform capable of supporting several complementary areas of astrophysics, with major programs in wide-field cosmology and galaxy evolution, direct imaging and characterization of nearby planetary systems, time-domain and multi-messenger observations across several programs, compact-object studies, and observations of Solar-System small bodies. These programs place common requirements on angular resolution, photometric stability, rapid target acquisition, spectroscopic capability, and long-term observing efficiency.
The telescope employs a segmented 3.5-meter primary mirror and is designed for high-angular-resolution imaging from the near-ultraviolet through the optical and near-infrared. The current baseline covers 0.2--1.5~$\mu$m, with the wavelength at which diffraction-limited performance is specified to be determined from the final wavefront-error budget.
A wide-field imaging system is intended to support deep surveys, precision photometry, and repeated monitoring over a field of view of approximately $10' \times 10'$ to $30' \times 30'$. Spectroscopic modes with resolving powers around $R \sim 1000$, together with higher-resolution options approaching $R \sim 5000$, are being considered for galaxy surveys, transient classification, compact-object spectroscopy, and other targeted investigations.
A dedicated coronagraphic capability is also being studied for direct observations of nearby exoplanetary systems. The current performance goal is a raw contrast of order $10^{-8}$, with further improvement expected through calibration and post-processing, while maintaining the angular resolution required to probe planetary systems around the nearest stars. Candidate mission configurations include operations near the Sun--Earth L2 region as well as alternative Earth-orbit architectures, with the final choice to be determined by science performance, thermal stability, communications, operations, and mission cost.
The telescope is intended to combine survey observations, long-duration monitoring, and time-sensitive observations. The telescope integrates wide-field imaging, spectroscopy, and 
coronagraphic capability around a common 18-segment optical system.
This paper defines the current science requirements, baseline technical
configuration, and engineering trade space for subsequent development
of the 3.5mST concept.
}

\newcommand{\wpkeywords}{\textbf{Keywords:} segmented mirror, robotic space telescope, cosmology, galaxy evolution, exoplanet, time-domain astronomy, multi-messenger astronomy}
